% =============================================================
% CAPÍTULO 6 — Síntesis: de la interpretabilidad a la compresión informada
% Versión: matrícula de honor
% =============================================================

\section{Síntesis: compresión informada por interpretabilidad}
\label{sec:sintesis}

Los dos capítulos anteriores han recorrido caminos complementarios. El
Capítulo~\ref{sec:resultados-compresion} midió qué se rompe al
comprimir BERT con SVD: la sensibilidad varía radicalmente entre
componentes (factor 14$\times$ a rango~128 en retención absoluta y hasta
72$\times$ a rango~64 normalizado por parámetros), entre profundidades
(transición de fase entre rango~384 y 256) y entre emociones (desde
\emph{gratitude}, que retiene el 99\% del F1 a rango~384, hasta
\emph{embarrassment}, que colapsa al 0\% al mismo rango). El
Capítulo~\ref{sec:resultados-interpretabilidad} investigó por qué
mediante cinco técnicas de interpretabilidad mecánica que convergen en una
arquitectura funcional de tres niveles: información léxica formada en
capas tempranas, cómputo de transición en capas medias, y alineación
geométrica con la base del clasificador en capas tardías, donde reside el
cuello de botella computacional del modelo.

Este capítulo cierra el bucle: utiliza los hallazgos de interpretabilidad
mecánica para diseñar estrategias de compresión \emph{informadas} y evalúa
si ese conocimiento se traduce en ventaja cuantitativa sobre las
estrategias ciegas presentadas en el
Capítulo~\ref{sec:resultados-compresion}. La estructura del capítulo
refleja una progresión metodológica (intuición, fracaso de la
intuición, pivot a datos, dominancia) que constituye en sí misma una
observación replicable sobre cómo trasladar conocimiento mecanicista a
ingeniería aplicada.

El recorrido se organiza en cinco secciones. La primera establece la
conexión cuantitativa entre cristalización y vulnerabilidad
(\S\ref{sec:cristalizacion-vulnerabilidad}). La segunda describe la
estrategia heurística informada y documenta su fracaso
(\S\ref{sec:informed-heuristic}). La tercera presenta el algoritmo greedy
alimentado con datos empíricos de sensibilidad y su dominancia de la
frontera de Pareto (\S\ref{sec:greedy-algorithm}). La cuarta evalúa la
recuperación mediante \emph{fine-tuning} y reporta una observación inesperada
sobre regularización implícita (\S\ref{sec:finetuning-recovery}). La quinta
sintetiza las implicaciones del conjunto de resultados
(\S\ref{sec:sintesis-discusion}).


% -------------------------------------------------------------
\subsection{La cristalización predice la vulnerabilidad}
\label{sec:cristalizacion-vulnerabilidad}
% -------------------------------------------------------------

La primera pregunta de síntesis es si los dos ejes del proyecto
(compresión e interpretabilidad) convergen cuantitativamente, no
solo metafóricamente. El Capítulo~\ref{sec:resultados-interpretabilidad}
mostró que cada emoción cristaliza a una profundidad distinta:
\emph{gratitude} alcanza el 80\% de su F1 máximo en la capa~0, mientras que
\emph{disappointment} no lo hace hasta la capa~9. El
Capítulo~\ref{sec:resultados-compresion} mostró que la compresión SVD
uniforme afecta a cada emoción de forma diferencial.

\subsubsection{Correlación cristalización-vulnerabilidad}

La correlación entre la capa de cristalización y la caída de F1 bajo
compresión SVD uniforme a rango~128 es positiva y significativa (Spearman
$\rho \approx 0{,}48$, $p < 0{,}05$, $n=23$): las emociones que cristalizan
en capas más profundas son sistemáticamente más vulnerables a la
compresión. Este resultado es coherente con la estructura espectral
documentada en \S\ref{sec:analisis-espectral}: las matrices de las capas
tardías presentan distribuciones de valores singulares menos concentradas
(menor redundancia), lo que las hace más frágiles bajo truncamiento SVD.

La conexión, sin embargo, no se agota en una correlación. Tres predictores
independientes de la vulnerabilidad emergen del análisis combinado:

\begin{itemize}[nosep]
    \item Profundidad de cristalización por probing (\S\ref{sec:probing}):
    $\rho \approx 0{,}48$.
    \item Aislamiento representacional, distancia media a otras emociones
    en el espacio de \texttt{[CLS]} de L11 (\S\ref{sec:probing}):
    $\rho = 0{,}45$, $p = 0{,}031$.
    \item Norma del vector de selectividad neuronal
    (\S\ref{sec:neuron-analysis}): $\rho = 0{,}64$, $p = 0{,}001$.
\end{itemize}

El predictor más fuerte es la \emph{intensidad de codificación neuronal}
, no la geometría SVD ni el tamaño del dataset, ambos no significativos.
Este resultado es teóricamente importante: la SVD no ataca selectivamente
las neuronas emocionales (la fracción de señal en el subespacio eliminado
es prácticamente idéntica para las 23 emociones, rango 0{,}53--0{,}56),
pero las emociones que el modelo escribe con mayor intensidad neuronal son
más sensibles a perturbaciones generales de los pesos. Es una
vulnerabilidad por intensidad, no por geometría.

La Figura~\ref{fig:cap6-emotional-landscape} sintetiza visualmente esta
estructura cruzando los dos predictores principales: la capa de
cristalización (eje horizontal) y la norma de selectividad neuronal
(eje vertical).
 
\begin{figure}[H]
    \centering
    \includegraphics[width=0.95\textwidth]{figures/cap6/cap6_emotional_landscape_es.png}
    \caption{Mapa emocional del modelo. Cada punto representa una de las
    23 emociones, situada según su capa de cristalización por probing
    (eje horizontal, derivado del Capítulo~\ref{sec:resultados-interpretabilidad})
    y su norma de selectividad neuronal (eje vertical, intensidad total
    con la que el modelo codifica esa emoción a través de sus neuronas
    FFN). Los colores agrupan las emociones por cluster semántico. Las
    bandas verticales sombreadas marcan los tres bloques de profundidad:
    temprano (Emb--L3), medio (L4--L7) y tardío (L8--L11). Las emociones
    léxicas (\emph{gratitude}, \emph{love}, \emph{amusement}) ocupan la
    región superior izquierda; las contextuales (\emph{annoyance},
    \emph{disappointment}) la inferior derecha.}
    \label{fig:cap6-emotional-landscape}
\end{figure}
 
La distribución no es uniforme. Las emociones léxicas se concentran en
la región superior izquierda del mapa: cristalización temprana y alta
intensidad neuronal. \emph{Gratitude} es el caso extremo, con
cristalización en L0 y la mayor intensidad de codificación neuronal del
dataset. Las emociones contextuales ocupan la región inferior derecha:
cristalización tardía y baja intensidad. \emph{Disappointment} y
\emph{annoyance} muestran una intensidad de codificación
aproximadamente un orden de magnitud inferior a pesar de cristalizar
en L9--L10. Los valores absolutos del eje vertical de la
Figura~\ref{fig:cap6-emotional-landscape} corresponden a la \emph{máxima}
norma capa a capa del vector de selectividad (resumen visual), mientras
que la norma agregada $\ell_2$ a través de las 12~capas, reportada en
\S\ref{sec:neuron-analysis} y utilizada como predictor cuantitativo
($\rho = 0{,}639$), produce la misma ordenación entre emociones pero
en una escala distinta (\emph{gratitude}: 140; \emph{remorse}: 111;
\emph{disappointment}: 42).
 
Esta correspondencia entre dos magnitudes derivadas de experimentos
independientes ---probing por un lado, análisis de selectividad por
otro--- no es trivial. Probing mide la separabilidad lineal de la
representación; selectividad mide la especialización de las neuronas
individuales. Que ambas correlacionen sugiere que la facilidad de
aprendizaje y la intensidad de codificación son manifestaciones de una
misma propiedad subyacente: la disponibilidad de señal léxica robusta.
Las emociones contextuales no son simplemente más difíciles ---son
representacionalmente más débiles---, y esa debilidad estructural es
la que las hace particularmente vulnerables a la compresión.

\subsubsection{Tres principios de compresión informada}

De estos hallazgos se derivan tres principios que orientan el diseño de
estrategias de compresión informadas:

\begin{enumerate}
    \item Proteger las capas tardías. Concentran la
    especialización funcional (84\% de neuronas significativas, 77\% de
    cabezas críticas) y tienen menor redundancia espectral. Comprimir
    capas 8--11 al mismo rango que capas 0--3 es estructuralmente
    incorrecto.

    \item Distinguir componentes dentro de cada capa. La FFN
    Intermediate es el cuello de botella computacional confirmado por
    \emph{activation patching} (\S\ref{sec:activation-patching}); Q y K toleran
    compresión agresiva con coste despreciable. Aplicar el mismo rango a
    Q y a FFN Intermediate ignora un factor 14$\times$ de diferencia en
    retención de F1 a rango~128 (que se amplía a 72$\times$ a rango~64
    normalizado por parámetros eliminados).

    \item Aceptar que la compresión introduce sesgo emocional.
    Las emociones contextualmente complejas (cristalización tardía,
    codificación intensa) son penalizadas desproporcionadamente. Este
    sesgo es \emph{inevitable} sin intervención selectiva: cualquier
    estrategia uniforme lo introduce de forma silenciosa.
\end{enumerate}

Los principios son cualitativos. La pregunta abierta, y que motiva el
resto del capítulo, es si pueden traducirse directamente en reglas
heurísticas o si requieren cuantificación empírica.


% -------------------------------------------------------------
\subsection{Estrategia heurística: un resultado negativo informativo}
\label{sec:informed-heuristic}
% -------------------------------------------------------------

El primer intento de traducir la comprensión mecánica en una estrategia de
compresión adoptó un enfoque heurístico: asignar rangos SVD a cada
componente y capa según reglas cualitativas derivadas de los hallazgos de
interpretabilidad. La premisa de partida era razonable: si comprender el
modelo permite comprimirlo mejor, las reglas explícitas derivadas de esa
comprensión deberían superar a las estrategias ciegas.

\subsubsection{Diseño de las tres variantes}

Se diseñaron tres niveles de agresividad respetando los principios de
\S\ref{sec:cristalizacion-vulnerabilidad}:

\begin{itemize}[nosep]
    \item \textbf{Informed-light}: compresión conservadora que protege las
    capas tardías con rangos altos (700 en capa~11) y aplica reducciones
    moderadas en capas tempranas. Ratio de parámetros: 0{,}97.

    \item \textbf{Informed-moderate}: compresión intermedia con protección
    de la capa~11 (rango 512) y rangos entre 170 y 430 en el resto. Ratio:
    0{,}72.

    \item \textbf{Informed-aggressive}: compresión fuerte con rangos entre
    86 y 384. Ratio: 0{,}47.
\end{itemize}

Las tres variantes respetaban estrictamente el gradiente de criticidad por
profundidad (más rango en capas tardías que en tempranas) y la asimetría
entre componentes (más rango en FFN que en atención). La intuición
mecanicista estaba codificada explícitamente en las reglas.

\subsubsection{Resultados}

Los resultados, sin embargo, son reveladores
(Tabla~\ref{tab:informed-heuristic}):

\begin{table}[ht]
\centering
\caption{Rendimiento de las estrategias heurísticas informadas frente a
baselines de complejidad comparable. Las tres variantes \emph{igualan
exactamente} al baseline uniforme con el mismo ratio de parámetros
(\texttt{informed\_aggressive} = \texttt{uniform\_r256};
\texttt{informed\_moderate} = \texttt{uniform\_r512}), o requieren más
parámetros que el original sin alcanzar la retención de
\texttt{adaptive\_e99} (\texttt{informed\_light}). Ninguna domina la
frontera de Pareto.}
\label{tab:informed-heuristic}
\small
\begin{tabular}{llccc}
\toprule
\textbf{Estrategia} & \textbf{Tipo} & \textbf{Ratio params.} & \textbf{F1 macro} & \textbf{Retención} \\
\midrule
uniform\_r256        & uniforme   & 0{,}612 & 0{,}025 & 4{,}3\% \\
informed\_aggressive & informada  & 0{,}612 & 0{,}025 & 4{,}3\% \\
\midrule
adaptive\_e80        & adaptativa & 0{,}721 & 0{,}143 & 24{,}8\% \\
\midrule
uniform\_r512        & uniforme   & 1{,}000 & 0{,}464 & 80{,}4\% \\
informed\_moderate   & informada  & 1{,}000 & 0{,}464 & 80{,}4\% \\
\midrule
adaptive\_e95        & adaptativa & 1{,}022 & 0{,}491 & 85{,}0\% \\
adaptive\_e99        & adaptativa & 1{,}202 & 0{,}558 & 96{,}7\% \\
informed\_light      & informada  & 1{,}285 & 0{,}561 & 97{,}2\% \\
\bottomrule
\end{tabular}
\end{table}

Las tres variantes presentan el mismo patrón estructural: \emph{igualan
exactamente} al baseline uniforme de complejidad equivalente, o lo superan
solo a costa de un incremento sustancial de parámetros. La variante
agresiva, con ratio 0{,}612 y F1 = 0{,}025, reproduce número a número los
resultados de \texttt{uniform\_r256}: misma asignación efectiva de rangos
agregada, mismo colapso del modelo. La moderada, con ratio 1{,}000 y
F1 = 0{,}464, iguala a \texttt{uniform\_r512} con idéntica precisión.
La variante conservadora alcanza un F1 de 0{,}561 ---ligeramente
superior al baseline \texttt{adaptive\_e99}--- pero requiere un 28{,}5\%
\emph{más} parámetros que el modelo original (ratio 1{,}285), lo que la
sitúa fuera del régimen de compresión efectiva.

Ninguna de las tres variantes informadas aparece en la frontera de Pareto:
todas son \emph{dominadas} por las baselines de complejidad equivalente o
por las estrategias adaptativas energéticas.

\subsubsection{Diagnóstico del fracaso}

Este resultado negativo es \emph{instructivo}, no anecdótico. La intuición
cualitativa derivada de la interpretabilidad (``proteger capas tardías,
comprimir agresivamente Q y K, respetar la asimetría componente'') es
correcta en su dirección ---el algoritmo greedy
(\S\ref{sec:greedy-algorithm}) confirma exactamente esa secuencia de
decisiones cuando opera con datos cuantitativos---, pero traducirla a
reglas heurísticas explícitas no produce ganancia alguna sobre las
baselines: las informed simplemente reproducen la misma frontera que las
estrategias ciegas alcanzan sin invocar interpretabilidad.

La razón es la \emph{magnitud relativa} de las sensibilidades. Saber que
Q es más compresible que FFN Intermediate no indica cuánto más, ni en
qué profundidad la asimetría es mayor, ni cuál es el orden exacto de
prioridad entre Q en capas tardías vs.\ FFN Output en capas tempranas. La
regla heurística asigna, por ejemplo, ``rango 256 a Q en capas medias'';
los datos cuantitativos revelan que Q a rango 128 retiene 99{,}4\% del
F1 (\S\ref{sec:sensibilidad-componente}), por lo que el rango 256 es
innecesariamente conservador. Multiplicado por las 72 matrices del modelo,
el conservadurismo acumulado lleva a la heurística a converger
exactamente sobre los puntos que la compresión uniforme alcanza sin
ninguna información estructural: \texttt{informed\_aggressive} reproduce
\texttt{uniform\_r256}, \texttt{informed\_moderate} reproduce
\texttt{uniform\_r512}. La heurística \emph{no extrae valor adicional}
del conocimiento mecanicista que pretende traducir.

La conclusión metodológica es precisa: \emph{traducir interpretabilidad
mecánica en compresión efectiva requiere datos empíricos de sensibilidad,
no solo comprensión arquitectónica}. La interpretabilidad cualitativa
identifica qué investigar; la sensibilidad cuantitativa determina cuánto
comprimir. Que las tres variantes informadas converjan sobre la frontera
de las baselines ciegas, en lugar de superarla, es la evidencia más
limpia posible de esta separación: el conocimiento cualitativo de la
estructura funcional es necesario pero no suficiente para guiar la
compresión. Esta distinción motiva el cambio de enfoque hacia el
algoritmo greedy data-driven.


% -------------------------------------------------------------
\subsection{Algoritmo greedy: compresión guiada por datos}
\label{sec:greedy-algorithm}
% -------------------------------------------------------------

El fracaso de la estrategia heurística motivó un cambio de enfoque: en
lugar de traducir la intuición mecánica en reglas, utilizar directamente
los datos cuantitativos de sensibilidad por componente
(\S\ref{sec:sensibilidad-componente}) y por profundidad como entrada para
un algoritmo de optimización.

\subsubsection{Por qué greedy y no otra estrategia}

El problema es seleccionar simultáneamente 18~rangos (6~componentes
$\times$ 3~bandas de profundidad) en un menú discreto de 3~valores
candidatos cada uno, optimizando el F1 macro bajo una restricción de
parámetros máximos. Tres aproximaciones aparecen como candidatas
naturales y todas presentan inconvenientes serios.

La búsqueda exhaustiva es inviable: $3^{18} \approx 3{,}9\cdot 10^8$
configuraciones, $\sim$30~s por evaluación en GPU T4, dan unos
370~años de cómputo. La optimización bayesiana sobre el espacio de
rangos discreto es viable conceptualmente, pero exige cientos de
evaluaciones del modelo comprimido para construir el \emph{surrogate}
(varias horas de GPU) y el espacio discreto degrada la calidad del GP.
Los métodos diferenciables (\textit{straight-through estimators},
Gumbel-softmax sobre rangos) requieren relajaciones que cambian la
naturaleza del problema y no aprovechan los datos de sensibilidad ya
medidos.

El greedy data-driven evita estas tres limitaciones. Es interpretable
---la secuencia de movimientos se puede leer y se corresponde con una
narrativa funcional que valida cruzadamente los hallazgos del
Capítulo~\ref{sec:resultados-interpretabilidad}---, lineal en el
número de movimientos ($O(|moves|)$) y resuelve el problema en menos
de un segundo, sin requerir ninguna evaluación adicional del modelo
porque los costes provienen de los datos pre-calculados de
sensibilidad. Como contrapartida es sub-óptimo: el algoritmo no
explora el espacio combinatorio completo y no garantiza alcanzar el
máximo global. Su validación es empírica: 8~de los 9~puntos
Pareto-óptimos sobre 21~configuraciones evaluadas son greedy
(\S\ref{sec:greedy-algorithm}), lo que sugiere que la sub-optimalidad
real es modesta dentro de la región funcional del espacio
compresión--rendimiento.

\subsubsection{Diseño del algoritmo}

El algoritmo greedy opera sobre un menú finito de \emph{movimientos de
compresión}, donde cada movimiento consiste en reducir el rango SVD de un
componente (Q, K, V, Attn Output, FFN Intermediate, FFN Output) en una
banda de profundidad (Early: capas 0--3, Middle: 4--7, Late: 8--11) desde
su rango actual a un rango inferior. Los rangos candidatos no son
uniformes para todos los componentes: se restringen a aquellos donde la
pérdida de F1 medida en \S\ref{sec:sensibilidad-componente} es asumible.
En concreto, Q y K admiten los tres rangos $\{256, 128, 64\}$; V, Attn
Output y FFN Output sólo $\{256, 128\}$; FFN Intermediate sólo $\{256\}$,
porque a rango~128 ya colapsa al 6{,}9\% de retención y a rango~64 al
0\%. Esta restricción es un \emph{filtro previo} derivado de los
hallazgos de sensibilidad y reduce el menú a 14~movimientos por banda
($\times$\,3 bandas = 42 movimientos posibles). Para llegar al rango más
agresivo el algoritmo debe pasar por los intermedios: los movimientos son
incrementales.

Cada movimiento tiene dos magnitudes asociadas:

\begin{itemize}
    \item \textbf{Coste incremental} ($\Delta$F1): la caída de F1 macro
    al bajar del rango previo al actual (no respecto al baseline), de
    modo que el algoritmo trata la compresión como una secuencia de
    decisiones independientes. El coste se modula por la banda de
    profundidad según las proporciones empíricas de impacto observadas
    en \S\ref{sec:sensibilidad-componente}: las capas tempranas absorben
    el 14{,}8\% del coste total, las medias el 17{,}7\% y las tardías el
    67{,}5\%. Como filtro adicional, se descartan los movimientos cuyo
    coste incremental excede 15~puntos porcentuales (\textit{cap} de
    seguridad \texttt{max\_cost\_per\_move=15.0} en la implementación),
    lo que evita decisiones catastróficas como comprimir FFN Intermediate
    en capas tardías.

    \item \textbf{Ahorro}: la reducción en número de parámetros al bajar el
    rango. Un detalle técnico crítico es el \emph{rango de equilibrio}
    (\textit{break-even rank}): la descomposición SVD $W \approx U \Sigma V^T$
    solo reduce parámetros cuando $r < mn/(m+n)$. Para las matrices de
    atención ($768 \times 768$), el equilibrio está en rango~384; para las
    matrices FFN ($768 \times 3072$), en rango~614{,}4. Rangos superiores a
    estos umbrales \emph{incrementan} el número de parámetros, lo que el
    algoritmo evita por construcción.
\end{itemize}

La eficiencia de cada movimiento se define como el cociente entre el
ahorro de parámetros y el coste incremental en F1:
\begin{equation}
    \mathcal{E}(\text{movimiento}) =
    \begin{cases}
        +\infty & \text{si } \Delta F1 \leq 0 \\[4pt]
        \dfrac{\text{params\_ahorrados}}{\Delta F1 \cdot 100} & \text{en otro caso}
    \end{cases}
    \label{eq:eficiencia}
\end{equation}
El algoritmo selecciona movimientos en orden decreciente de eficiencia
hasta alcanzar un ratio de parámetros objetivo. Los movimientos con coste
incremental nulo o negativo (componentes que mejoran ligeramente al
comprimirse, lo que ocurre, por ejemplo, con Q a rango~256 con retención
del 100{,}4\%) tienen eficiencia infinita y se aplican primero. La
formalización completa se recoge en el Algoritmo~\ref{alg:greedy}.

\begin{algorithm}[H]
\caption{Greedy SVD Compression (informed by sensitivity)}
\label{alg:greedy}
\begin{algorithmic}[1]
\Require Modelo $M$, ratio objetivo de parámetros $r_{\text{tgt}}$,
         tabla de retención $S[c][k]$, fracciones de profundidad $f[d]$,
         cap de coste $\tau$
\Ensure Asignación de rangos $\sigma:(c,d)\!\to\!\text{rank}\cup\{\text{full}\}$
\State $\text{moves}\gets\emptyset$
\For{cada componente $c$}
  \For{cada par consecutivo $(k_{\text{prev}},k)$ en $S[c]$}
    \For{cada banda de profundidad $d\in\{\text{early},\text{mid},\text{late}\}$}
      \State $\Delta F1 \gets \big(S[c][k_{\text{prev}}]-S[c][k]\big)\cdot f[d]$
      \State $\text{saved}\gets |d|\cdot\big(k_{\text{prev}}(m{+}n)-k(m{+}n)\big)$
      \State $\text{eff} \gets \begin{cases}+\infty & \Delta F1\le 0\\ \text{saved}/\Delta F1 & \text{otro caso}\end{cases}$
      \State $\text{moves}.\text{add}(\langle c,d,k_{\text{prev}}{\to}k,\Delta F1,\text{saved},\text{eff}\rangle)$
    \EndFor
  \EndFor
\EndFor
\State ordenar $\text{moves}$ por $\text{eff}$ descendente
\State $\sigma[c,d]\gets\text{full}$ para todo $(c,d)$
\For{$m\in\text{moves}$}
  \If{$\sigma[m.c,m.d]\ne m.k_{\text{prev}}$} \textbf{continue} \EndIf
  \If{$m.\Delta F1>\tau$} \textbf{continue} \EndIf
  \State $\sigma[m.c,m.d]\gets m.k$
  \If{$\text{params\_ratio}(\sigma)\le r_{\text{tgt}}$} \textbf{break} \EndIf
\EndFor
\State \Return $\sigma$
\end{algorithmic}
\end{algorithm}


\subsubsection{Secuencia de decisiones}

La secuencia con que el algoritmo asigna compresión revela una jerarquía
de importancia que valida \emph{independientemente} los hallazgos de
interpretabilidad (Tabla~\ref{tab:greedy-sequence}):

\begin{table}[ht]
\centering
\caption{Secuencia de decisiones del algoritmo greedy. Cada fila muestra
los componentes comprimidos al alcanzar el ratio indicado. El algoritmo
opera exclusivamente con datos numéricos de sensibilidad, sin acceso a
los hallazgos cualitativos de interpretabilidad.}
\label{tab:greedy-sequence}
\small
\begin{tabular}{ccp{8.0cm}}
\toprule
\textbf{Config.} & \textbf{Ratio} & \textbf{Componentes comprimidos} \\
\midrule
greedy\_95 & 0{,}950 & Solo Q y K \\
greedy\_90 & 0{,}864 & Q, K, FFN Output (capas tempranas) \\
greedy\_85 & 0{,}813 & Q, K, FFN Output (tempranas + medias) \\
greedy\_80 & 0{,}799 & Q, K, V (tempranas), FFN Output \\
greedy\_75 & 0{,}741 & Q, K, V, FFN Inter+Out (tempranas) \\
greedy\_70 & 0{,}670 & Todos los componentes, mayoritariamente
                       tempranos y medios \\
\bottomrule
\end{tabular}
\end{table}

Tres observaciones emergen de esta secuencia, cada una con su contraparte
en interpretabilidad:

\begin{enumerate}
    \item Q y K son esencialmente gratuitas de comprimir. El
    primer 5\% de reducción de parámetros proviene íntegramente de
    comprimir las proyecciones de Query y Key, con un coste en F1 inferior
    al 0{,}5\%. Esto es coherente con el análisis de sensibilidad por
    componente (\S\ref{sec:sensibilidad-componente}), que documentó
    retenciones del 99{,}4\% para Q y 98{,}4\% para K a rango~128, y con
    el solapamiento espectral elevado entre subespacios Q-K
    (hasta 0{,}63 para los 32 vectores singulares principales,
    \S\ref{sec:analisis-espectral}). Q y K codifican esencialmente
    redundancia, no información distintiva.

    \item La FFN Intermediate no se comprime en capas tardías.
    Ninguna configuración greedy con retención superior al 30\% toca la
    FFN Intermediate de las capas 8--11. El algoritmo descubre de forma
    independiente el cuello de botella identificado por activation
    patching (\S\ref{sec:activation-patching}): restaurar únicamente la
    FFN de la capa~11 recuperaba el 100\% del rendimiento del modelo. La
    convergencia entre dos métodos completamente disjuntos (uno mide
    sensibilidad puntual, otro mide suficiencia causal) en el mismo
    componente como \emph{único} crítico es una validación cruzada fuerte.

    \item Las capas tardías se protegen sistemáticamente. Incluso
    en las configuraciones más agresivas (\texttt{greedy\_70}, ratio
    0{,}67), las capas 8--11 retienen rangos sustancialmente mayores que
    las tempranas. Esto valida el gradiente de criticidad documentado por
    la ablación de cabezas (\S\ref{sec:head-ablation}): el 77\% de las
    cabezas en capas 8--11 son críticas frente al 27\% en capas 0--4.
\end{enumerate}

La convergencia entre la secuencia de decisiones del algoritmo (que opera
solo con datos numéricos de sensibilidad por componente y profundidad) y
los hallazgos cualitativos de interpretabilidad (que explican \emph{por
qué} esos componentes son sensibles) constituye una validación cruzada de
ambas líneas de análisis. Si interpretabilidad y compresión llegan a las
mismas conclusiones por caminos disjuntos, la probabilidad de que ambas
estén sistemáticamente sesgadas en la misma dirección es baja.


\subsubsection{Dominancia en la frontera de Pareto}

La Figura~\ref{fig:pareto-all-strategies} presenta la comparación completa
de las 21 configuraciones evaluadas (6 uniformes, 4 adaptativas, 3
heurísticas informadas y 8 greedy) en el espacio ratio de parámetros
vs.\ F1 macro.

\begin{figure}[ht]
\centering
\includegraphics[width=0.95\textwidth]{figures/cap6/cap6_pareto_all_strategies_es.png}
\caption{Frontera de Pareto para las cuatro familias de estrategias de
compresión. Los puntos sobre la curva envolvente representan
configuraciones Pareto-óptimas. 8 de los 9 puntos Pareto-óptimos
corresponden a estrategias greedy. La única estrategia no-greedy en la
frontera es \texttt{uniform\_r32}, que aparece solo por tener el menor
ratio absoluto pero con F1 = 0 (el modelo está colapsado).}
\label{fig:pareto-all-strategies}
\end{figure}

De los 9 puntos Pareto-óptimos en sentido estricto, 8 son estrategias
greedy. El algoritmo basado en datos empíricos domina la frontera
completa en la región funcional ($F1 > 0$), desde un ratio de
0{,}58 hasta 0{,}95. La Tabla~\ref{tab:head-to-head} cuantifica las
ventajas en comparaciones directas a ratios de parámetros similares:

\begin{table}[ht]
\centering
\caption{Comparaciones directas entre estrategias greedy y alternativas a
ratios de parámetros comparables. Los intervalos de confianza al 95\% se
obtienen mediante \emph{bootstrap por emoción} (10\,000 réplicas,
$n=23$); el \emph{p}-valor en la columna $\Delta$F1 corresponde a un
test \emph{paired bootstrap} sobre la diferencia por emoción cuando se
dispone de datos por emoción para ambas estrategias.}
\label{tab:head-to-head}
\small
\begin{tabular}{lcccc}
\toprule
\textbf{Comparación a ratio similar} & \textbf{Greedy F1 [CI95]} & \textbf{Alt.\ F1 [CI95]} & $\Delta$\textbf{F1} & \textbf{Sig.} \\
\midrule
greedy\_80 ($\rho=0{,}80$) vs uniform\_r384 ($\rho=0{,}81$)
  & 0{,}502\textsuperscript{$\dagger$}
  & 0{,}251 [0{,}151; 0{,}365]
  & $+$0{,}251 & --- \\
greedy\_90 ($\rho=0{,}86$) vs adaptive\_e90 ($\rho=0{,}89$)
  & 0{,}539 [0{,}467; 0{,}613]
  & 0{,}357 [0{,}255; 0{,}466]
  & $+$0{,}181 [$+$0{,}130; $+$0{,}241]
  & $p<0{,}001$ \\
greedy\_95 ($\rho=0{,}95$) vs adaptive\_e95 ($\rho=1{,}02$)
  & 0{,}575\textsuperscript{$\dagger$}
  & 0{,}491 [0{,}400; 0{,}581]
  & $+$0{,}084 & --- \\
\bottomrule
\multicolumn{5}{@{}l}{\footnotesize \textsuperscript{$\dagger$}Se reporta
únicamente el F1 macro puntual: los datos por emoción para
\texttt{greedy\_80} y \texttt{greedy\_95} no se exportaron a CSV en NB09. El
test \emph{paired bootstrap} para greedy\_90 vs adaptive\_e90 produce un
$\Delta$F1 = $+$0{,}181 con $p<0{,}001$ y CI estrictamente positivo, lo}\\
\multicolumn{5}{@{}l}{\footnotesize que es la evidencia estadística más
limpia disponible. Para los otros dos comparativos, los puntos del greedy
caen \emph{fuera} del intervalo de confianza al 95\% de la alternativa, lo
que es consistente con significancia direccional.}
\end{tabular}
\end{table}

El resultado más llamativo se da en el rango intermedio: a un ratio de
$\sim$0{,}80, la estrategia greedy retiene el 87{,}0\% del F1 original,
mientras que la uniforme a rango~384 solo retiene el 43{,}5\%: una
ventaja superior al doble en retención. La comparación
\texttt{greedy\_90} vs \texttt{adaptive\_e90}, donde sí se dispone de
datos por emoción para ambas, permite cuantificar formalmente la ventaja:
$\Delta$F1 = $+$0{,}181 con intervalo de confianza al 95\%
$[+0{,}130; +0{,}241]$ (\emph{paired bootstrap} sobre las 23 emociones,
10\,000 réplicas, $p<0{,}001$). El intervalo está estrictamente alejado
de cero, descartando la posibilidad de que la ventaja del greedy sea
atribuible a variabilidad muestral entre emociones. Nótese que la
comparación con \texttt{adaptive\_e95} es particularmente reveladora: el
algoritmo greedy obtiene un F1 superior con un 7\% \emph{menos} de
parámetros, dado que la estrategia adaptativa al 95\% de energía produce
un ratio de 1{,}02 (es decir, \emph{incrementa} los parámetros, como se
discutió en \S\ref{sec:compresion-adaptativa}).

A ratio 0{,}95, la estrategia \texttt{greedy\_95} alcanza una retención
del 99{,}6\% del F1 baseline (sin pérdida medible) con
una reducción del 5\% en parámetros que viene íntegramente de comprimir Q
y K, lo que ilustra que existe un margen de compresión genuinamente
gratuito que las estrategias ciegas son incapaces de explotar.

\subsubsection{Comparación cualitativa con destilación}
\label{sec:greedy-vs-distilbert}

Las cifras del greedy se entienden mejor situándolas frente a la
familia ortogonal de compresión más popular: la destilación de
conocimiento.
DistilBERT~\cite{sanh2019distilbert}, la destilación canónica de
BERT-base, reduce el modelo a 6~capas (66~M parámetros, ratio
0{,}60 sobre los 110~M de BERT-base) reteniendo, según los autores,
el 97\% del rendimiento de BERT-base en GLUE. Esta comparación no es
estrictamente trasladable a GoEmotions multi-etiqueta sin re-evaluar
DistilBERT en la misma tarea ---ejercicio que se propone como
trabajo futuro (\S\ref{sec:trabajo-futuro})---, pero permite tres
observaciones de calibración.

A nivel de compresión agresiva ($\rho \approx 0{,}60$), DistilBERT
domina cómodamente: nuestra \texttt{greedy\_60} retiene apenas el
5{,}5\% del F1, mientras que la destilación re-entrena el modelo
pequeño desde cero con supervisión de un \emph{teacher} y puede
reorganizar libremente sus pesos. La SVD opera sobre los pesos
existentes sin re-entrenamiento, una restricción que pasa factura en
el régimen de alta reducción.

En la zona de compresión moderada ($\rho \in [0{,}80, 0{,}90]$) la
situación se invierte. La SVD greedy es competitiva o superior a las
técnicas ciegas convencionales sin necesidad de re-entrenar. Y si se
admite \emph{fine-tuning} posterior, la combinación greedy + FT
(\S\ref{sec:finetuning-recovery}) supera incluso al baseline sin
comprimir, algo que ninguna técnica ciega consigue a este ratio.

Las dos familias son complementarias más que competidoras. La
destilación es la opción correcta cuando se tolera el coste de
pre-procesamiento y la latencia de desarrollo larga; la SVD informada
lo es cuando se requiere intervenir directamente sobre un modelo ya
desplegado, o cuando se buscan reducciones moderadas
($<25\%$ params) sin tocar el ciclo de entrenamiento. La combinación
de ambas ---destilación seguida de SVD sobre el destilado--- está
poco explorada y constituye un experimento natural derivado de este
trabajo.

La elección del SVD selectivo en este TFG no pretende superar a la
destilación en su propio terreno, sino habilitar una zona del espacio
compresión--rendimiento que la destilación deja sin cubrir: ratios
0{,}80--0{,}95 con intervención post-hoc sobre el modelo
ya \emph{fine-tuneado}.


% -------------------------------------------------------------
\subsection{Recuperación mediante \emph{fine-tuning}}
\label{sec:finetuning-recovery}
% -------------------------------------------------------------

La compresión SVD, incluso con asignación de rangos óptima, introduce una
perturbación en los pesos del modelo. Una pregunta natural es si un breve
re-entrenamiento puede recuperar, o incluso superar, el rendimiento
original.

Se seleccionó la configuración \texttt{greedy\_90} (86{,}4\% de los
parámetros, F1 macro = 0{,}539) como candidata para \emph{fine-tuning}, por
ofrecer un equilibrio atractivo entre compresión y rendimiento. El
re-entrenamiento consistió en 3~épocas sobre GoEmotions con los
hiperparámetros estándar (tasa de aprendizaje $2 \times 10^{-5}$,
\textit{warmup} 10\%, \textit{batch size} 32, idénticos a los del
\emph{fine-tuning} original).

\subsubsection{Resultado agregado}

\begin{table}[ht]
\centering
\caption{Rendimiento del modelo comprimido (\texttt{greedy\_90}) antes y
después de \emph{fine-tuning}, comparado con el baseline sin comprimir.}
\label{tab:finetuning-summary}
\small
\begin{tabular}{lccc}
\toprule
\textbf{Modelo} & \textbf{Parámetros} & \textbf{F1 macro} & \textbf{vs.\ baseline} \\
\midrule
Baseline (sin compresión) & 100{,}0\% & 0{,}577 & --- \\
greedy\_90 (comprimido)   & 86{,}4\%  & 0{,}539 & $-$6{,}7\% \\
greedy\_90 (\emph{fine-tuned})   & 86{,}4\%  & \textbf{0{,}591} & \textbf{$+$2{,}4\%} \\
\bottomrule
\end{tabular}
\end{table}

El modelo comprimido y re-entrenado supera al baseline original con un
13{,}6\% menos de parámetros. Tres épocas de \emph{fine-tuning} no solo recuperan
la caída de 6{,}7 puntos porcentuales, sino que añaden 1{,}4 puntos
adicionales (todos los porcentajes calculados sobre F1 macro = 0{,}577 del
baseline). Este resultado, observado consistentemente al repetir el
experimento con distintas semillas aleatorias, sugiere que la compresión
SVD, al eliminar componentes de bajo valor singular, actúa durante el
re-entrenamiento como una forma de regularización implícita.

Esta interpretación tiene una limitación metodológica: el experimento no incluye un grupo de control consistente
en el baseline original con tres épocas adicionales de \emph{fine-tuning}. Es
posible, aunque improbable dada la saturación de F1 macro observada
durante las 4 épocas iniciales, que parte de la mejora provenga
simplemente de más entrenamiento y no del efecto regularizador de la
compresión. El experimento de control que verificaría la hipótesis se propone en el Capítulo~\ref{sec:conclusiones}.


\subsubsection{Análisis por emoción}

El patrón de recuperación no es uniforme: las emociones con F1 baseline
bajo son las que más se benefician de la compresión seguida de
\emph{fine-tuning}, y este patrón es sistemático
(Tabla~\ref{tab:finetuning-emotions}).

\begin{table}[ht]
\centering
\caption{Emociones con mayor ganancia y mayor pérdida tras compresión +
\emph{fine-tuning}, respecto al baseline original. La columna $\Delta$ es la
diferencia entre F1 \emph{fine-tuned} y F1 baseline.}
\label{tab:finetuning-emotions}
\small
\begin{tabular}{lcccc}
\toprule
\textbf{Emoción} & \textbf{F1 base} & \textbf{F1 comprim.} & \textbf{F1 \emph{fine-tuned}} & \textbf{$\Delta$ vs.\ base} \\
\midrule
\multicolumn{5}{l}{\emph{Mayores ganancias}} \\
embarrassment & 0{,}267 & 0{,}273 & 0{,}509 & \textbf{$+$0{,}242} \\
desire        & 0{,}535 & 0{,}554 & 0{,}621 & $+$0{,}086 \\
excitement    & 0{,}454 & 0{,}431 & 0{,}497 & $+$0{,}043 \\
realization   & 0{,}271 & 0{,}195 & 0{,}298 & $+$0{,}027 \\
disappointment & 0{,}275 & 0{,}178 & 0{,}301 & $+$0{,}026 \\
\midrule
\multicolumn{5}{l}{\emph{Mayores pérdidas}} \\
caring        & 0{,}542 & 0{,}479 & 0{,}504 & $-$0{,}038 \\
admiration    & 0{,}757 & 0{,}695 & 0{,}727 & $-$0{,}029 \\
fear          & 0{,}710 & 0{,}680 & 0{,}691 & $-$0{,}018 \\
confusion     & 0{,}515 & 0{,}443 & 0{,}497 & $-$0{,}018 \\
sadness       & 0{,}632 & 0{,}575 & 0{,}621 & $-$0{,}011 \\
\bottomrule
\end{tabular}
\end{table}

El caso de \emph{embarrassment} es el más extremo: el modelo comprimido
mantiene una capacidad similar al baseline (F1 = 0{,}273 frente a 0{,}267),
pero el \emph{fine-tuning} la eleva hasta F1 = 0{,}509, un 90{,}6\% superior al
baseline original. Una observación parecida ocurre con \emph{desire}: el
modelo comprimido apenas pierde rendimiento (0{,}554 vs 0{,}535), pero
tras tres épocas de \emph{fine-tuning} alcanza 0{,}621, un $+16\%$ relativo.
Este patrón, en el que la compresión apenas afecta y el \emph{fine-tuning}
amplifica, encaja con la interpretación de la SVD como
regularizador estructural: al eliminar direcciones de baja energía, la
compresión proporciona al re-entrenamiento un punto de partida sin
sobreajuste a patrones espurios. Las emociones infrarrepresentadas
(\emph{embarrassment}, \emph{desire}, \emph{realization},
\emph{disappointment}) son las que más se benefician de este efecto.

Las pérdidas, por contraste, son pequeñas (máximo 0{,}038 en
\emph{caring}) y afectan a emociones que ya tenían un F1 razonablemente
alto. La compresión seguida de \emph{fine-tuning} actúa como un ecualizador:
recorta ligeramente las emociones bien representadas y amplifica las
infrarrepresentadas. Este efecto es consistente con la eliminación de
componentes de alta varianza y baja contribución que caracteriza al
truncamiento SVD: el modelo, al perder direcciones que codificaban
sobreajuste a patrones frecuentes, recupera capacidad para las clases
minoritarias.


\subsubsection{Plasticidad asimétrica: qué cambia durante el \emph{fine-tuning}}
 
Una pregunta complementaria es \emph{dónde} reorganiza el modelo sus pesos
durante el \emph{fine-tuning} de recuperación. Si el efecto regularizador
hipotetizado en el apartado anterior es real, cabría esperar que el
re-entrenamiento concentrara su esfuerzo en los componentes funcionalmente
más relevantes para la tarea. La Figura~\ref{fig:cap6-finetuning-diff}
mide esta plasticidad como la magnitud relativa del cambio de cada
matriz, $\|\Delta W\| / \|W_0\|$, donde $\|\cdot\|$ es la norma de
Frobenius y $W_0$ el peso del modelo comprimido antes del \emph{fine-tuning}.
 
\begin{figure}[H]
    \centering
    \includegraphics[width=0.85\textwidth]{figures/cap6/cap6_finetuning_diff_es.png}
    \caption{Magnitud relativa del cambio de pesos durante las tres
    épocas de \emph{fine-tuning} posterior a la compresión, calculada como
    $\|\Delta W\| / \|W_0\|$ por capa (filas) y componente (columnas).
    La FFN de la capa~11 es el componente más reorganizado (FFN-i:
    5{,}1\%, FFN-o: 5{,}4\%), seguido de la atención de la misma capa
    (Attn-O: 4{,}6\%). Las proyecciones Q y K son las menos modificadas
    en todas las capas (3{,}0\%--3{,}4\%). El patrón refleja una
    asignación selectiva del esfuerzo de re-entrenamiento.}
    \label{fig:cap6-finetuning-diff}
\end{figure}
 
El resultado proporciona una validación cruzada con los hallazgos de
sensibilidad del Capítulo~\ref{sec:resultados-compresion}: los
componentes que más se reorganizan durante el \emph{fine-tuning} son
precisamente los más sensibles a la compresión. La FFN Intermediate de
la capa~11, cuyo colapso bajo compresión a rango 128 reducía el F1 al
7\% (\S\ref{sec:sensibilidad-componente}), es también el componente con
mayor cambio durante el re-entrenamiento. Las proyecciones Q y K, que
toleraban compresión a rango 128 con retención del 99\% del F1, son
también las que menos cambian al re-entrenarse.
 
Este paralelismo sugiere que sensibilidad a la compresión y plasticidad
bajo \emph{fine-tuning} miden, desde ángulos distintos, una misma propiedad
subyacente: la importancia funcional del componente para la tarea. Los
componentes que la tarea necesita más son los que más se reorganizan al
re-entrenarse, y los que más se rompen al comprimirse. Esta convergencia
añade una tercera evidencia ---junto a la sensibilidad por componente y
la convergencia del algoritmo greedy--- de que la jerarquía funcional
documentada por interpretabilidad mecánica no es un artefacto de un
método particular, sino una propiedad estructural del modelo.


% -------------------------------------------------------------
\subsection{Discusión y síntesis de hallazgos}
\label{sec:sintesis-discusion}
% -------------------------------------------------------------

Los resultados de este capítulo permiten extraer tres conclusiones
principales que integran las líneas de compresión e interpretabilidad del
proyecto.

\subsubsection{La interpretabilidad mecánica tiene valor aplicado, pero
requiere cuantificación}

El resultado negativo de las estrategias heurísticas
(\S\ref{sec:informed-heuristic}) y el éxito del algoritmo greedy
(\S\ref{sec:greedy-algorithm}) establecen un matiz importante para la
literatura de compresión informada: la interpretabilidad mecánica no
genera valor mediante la intuición cualitativa, sino a través de la
medición sistemática de sensibilidades. Las heurísticas informed
convergen exactamente sobre la frontera de las baselines ciegas
(\texttt{informed\_aggressive} $\equiv$ \texttt{uniform\_r256},
\texttt{informed\_moderate} $\equiv$ \texttt{uniform\_r512}), evidencia
limpia de que el conocimiento cualitativo no extrae ganancia adicional
sobre lo que la compresión uniforme alcanza sin invocar interpretabilidad.
Saber que ``la FFN es más importante que la atención'' no basta; lo que
permite diseñar estrategias superiores es cuantificar que Q a rango~256
retiene el 100{,}4\% del F1 mientras que FFN Intermediate a rango~64
colapsa al 0\%.

Esto no invalida la interpretabilidad cualitativa: los hallazgos de
probing, \emph{logit lens}, \emph{activation patching} y ablación siguen siendo
esenciales para explicar el comportamiento del modelo, pero
delimita su utilidad para la ingeniería de compresión: la interpretabilidad
identifica qué investigar, y la sensibilidad empírica determina
cuánto comprimir. La distinción es metodológicamente importante
porque permite predecir cuándo otros trabajos de ``compresión informada
por interpretabilidad'' tendrán éxito o fracasarán, en función de si se
basan en reglas heurísticas o en datos cuantitativos.

\subsubsection{Convergencia de cuatro métodos independientes}

La Tabla~\ref{tab:triangulation} resume la convergencia entre los
diferentes métodos experimentales del proyecto respecto a las tres
conclusiones estructurales principales:

\begin{table}[ht]
\centering
\caption{Triangulación de hallazgos: cuatro métodos independientes
convergen en las mismas conclusiones estructurales.}
\label{tab:triangulation}
\small
\begin{tabular}{p{3.0cm}p{3.3cm}p{3.5cm}p{3.5cm}}
\toprule
\textbf{Conclusión} & \textbf{Compresión} & \textbf{Interpretabilidad} & \textbf{Síntesis} \\
\midrule
Q/K son redundantes &
  Retención $>$99\% a rango 128 (\S\ref{sec:sensibilidad-componente}) &
  Solapamiento espectral 0{,}63 con K (\S\ref{sec:analisis-espectral});
  bajo impacto en ablación (\S\ref{sec:head-ablation}) &
  Primeros en comprimirse en greedy (\S\ref{sec:greedy-algorithm}) \\
\addlinespace
FFN es el cuello de botella &
  Asimetría 14$\times$ Q vs.\ FFN Intermediate (\S\ref{sec:sensibilidad-componente}) &
  Restaurar FFN L11 = 100\% (\S\ref{sec:activation-patching});
  84\% neuronas significativas en L8--11 (\S\ref{sec:neuron-analysis}) &
  Nunca comprimida en capas tardías (\S\ref{sec:greedy-algorithm}) \\
\addlinespace
Capas tardías son críticas &
  50{,}5\% del impacto total (\S\ref{sec:sensibilidad-componente}) &
  77\% cabezas críticas en L8--11 (\S\ref{sec:head-ablation});
  cristalización tardía predice F1 (\S\ref{sec:probing});
  cristalización geométrica L9--L11 (\S\ref{sec:logit-lens}) &
  Protegidas sistemáticamente (\S\ref{sec:greedy-algorithm}) \\
\bottomrule
\end{tabular}
\end{table}

Esta convergencia es significativa porque los métodos son
\emph{independientes en su diseño experimental}: el análisis de
sensibilidad mide F1 bajo perturbación de pesos; la ablación de cabezas
mide F1 bajo anulación de componentes; el probing mide separabilidad
lineal de representaciones; el \emph{logit lens} mide alineación con la base del
clasificador; y el \emph{activation patching} mide restauración funcional. Que
los cinco coincidan en las mismas conclusiones refuerza la validez de cada
método individual y, simultáneamente, la robustez de las conclusiones
estructurales. La probabilidad de que cinco métodos disjuntos compartan
exactamente el mismo sesgo sistemático es despreciable.

\subsubsection{La compresión como técnica de regularización: una hipótesis
con apoyo preliminar}

El resultado de \emph{fine-tuning} post-compresión (un modelo con 13{,}6\%
menos de parámetros que supera al original en F1 macro) invita a
reinterpretar la compresión SVD no solo como reducción de tamaño, sino
como una técnica de regularización implícita. El truncamiento de valores
singulares bajos elimina componentes de los pesos que contribuyen poco a
la varianza explicada pero pueden codificar ruido o sobreajuste. El
re-entrenamiento posterior permite al modelo reorganizar su
representación sobre una base más compacta.

Este efecto es particularmente visible en las emociones infrarrepresentadas
en el dataset (\emph{embarrassment}, \emph{realization}, \emph{desire}),
donde el baseline original sufre de sobreajuste a patrones espurios que la
compresión elimina. La mejora de 90{,}6\% en F1 de \emph{embarrassment}
(de 0{,}267 a 0{,}509) es el ejemplo más extremo, pero el patrón se
observa consistentemente en emociones de baja frecuencia.

La interpretación causal de este patrón (compresión $\rightarrow$
regularización $\rightarrow$ mejora) requiere experimentos de control
que no se han realizado en este trabajo y que se proponen en
\S\ref{sec:trabajo-futuro}. La observación que sí se sostiene es: \emph{la compresión greedy seguida de fine-tuning
produce un modelo simultáneamente más pequeño y más preciso que el
baseline, con ganancias concentradas en clases minoritarias}. Si la
hipótesis de regularización implícita se confirma con grupos de control
adecuados, la compresión SVD podría incorporarse como paso estándar en
pipelines de entrenamiento, no solo como técnica de reducción de tamaño.

La implicación práctica es clara independientemente de la causalidad
subyacente: un pipeline de \emph{fine-tuning $\rightarrow$ compresión
greedy $\rightarrow$ re-\emph{fine-tuning}} produce modelos que son
simultáneamente más pequeños y al menos tan buenos como el baseline
original, siempre que la compresión respete la estructura funcional del
modelo ---lo cual requiere, a su vez, el tipo de análisis de
interpretabilidad mecánica presentado en esta memoria.